% ===================================================================
%  estilo-apunte.tex  --  identidad visual FIQ-UNL para el apunte
%  de Análisis Matricial.
%
%  Se carga en main.tex con  \input{\VA estilo-apunte}  ANTES de las
%  declaraciones \newtheorem (redefine los estilos plain/definition/
%  remark, así que todos los entornos, incluidos los que se declaran
%  dentro de los capítulos, toman el estilo nuevo).
%
%  Marca FIQ: rojo #E42313, gris #585856, tipografía Lato.
%  Criterio: estructura en gris oscuro, rojo sólo como acento.
%  Lato para títulos, cajas y encabezados; el cuerpo en Latin Modern.
%
%  Paquetes que necesita (MiKTeX los instala solo): lmodern, lato,
%  microtype, geometry, titlesec, titletoc, fancyhdr, tcolorbox,
%  caption.  Los logos están en figuras/logo-fiq-*.png
% ===================================================================
\makeatletter

% ---------- tipografía ----------------------------------------------
\usepackage[T1]{fontenc}
\usepackage{lmodern}
\usepackage{lato}                 % da \lato (familia de la marca)
\usepackage{microtype}            % mejor justificado y \textls
\newcommand{\fiqfuente}{\lato}
\newcommand{\espaciado}[1]{\textls[160]{#1}}   % versalitas espaciadas

% ---------- página ---------------------------------------------------
\usepackage[a4paper,top=2.7cm,bottom=2.6cm,left=2.5cm,right=2.5cm,
  headheight=15pt,headsep=0.9cm,footskip=1.1cm]{geometry}

% ---------- colores de la marca -------------------------------------
\definecolor{fiqrojo}{HTML}{E42313}
\definecolor{fiqgris}{HTML}{585856}
\definecolor{fiqtexto}{HTML}{262625}     % casi negro, para títulos
\definecolor{fiqclaro}{HTML}{F5F5F4}     % fondo de cajas
\definecolor{fiqlinea}{HTML}{D6D6D3}     % filetes finos
\colorlet{fiqrojooscuro}{fiqrojo!78!black}

% ---------- enlaces --------------------------------------------------
\hypersetup{colorlinks=true,linkcolor=fiqrojooscuro,
  citecolor=fiqrojooscuro,urlcolor=fiqrojooscuro}

% ---------- títulos de capítulo y sección ---------------------------
\usepackage{titlesec}
\newcommand{\fiqcirculo}[2][1.55cm]{%
  \tikz[baseline=(c.base)]\node[circle,fill=fiqrojo,text=white,
    minimum size=#1,inner sep=0pt,font=\fiqfuente\bfseries\fontsize{24}{24}\selectfont](c){#2};}

\titleformat{\chapter}[display]
  {\fiqfuente\color{fiqtexto}}
  {\fiqcirculo{\thechapter}\hspace{0.9em}%
   {\small\color{fiqgris}\espaciado{\MakeUppercase{\chaptertitlename}}}}
  {16pt}
  {\fontsize{25}{30}\selectfont\bfseries\boldmath\raggedright}
  [\vspace{10pt}{\color{fiqlinea}\titlerule[0.6pt]}]
\titleformat{name=\chapter,numberless}[display]
  {\fiqfuente\color{fiqtexto}}
  {}
  {0pt}
  {\fontsize{25}{30}\selectfont\bfseries\boldmath\raggedright}
  [\vspace{10pt}{\color{fiqlinea}\titlerule[0.6pt]}]
\titlespacing*{\chapter}{0pt}{0pt}{28pt}

\titleformat{\section}
  {\fiqfuente\Large\bfseries\boldmath\color{fiqtexto}}
  {\textcolor{fiqrojo}{\thesection}}{0.7em}{}
\titleformat{name=\section,numberless}
  {\fiqfuente\Large\bfseries\boldmath\color{fiqtexto}}
  {}{0pt}{\textcolor{fiqrojo}{\rule[0.08em]{0.55em}{0.55em}}\hspace{0.5em}}
\titleformat{\subsection}
  {\fiqfuente\large\bfseries\boldmath\color{fiqtexto}}
  {\textcolor{fiqrojo}{\thesubsection}}{0.7em}{}
\titleformat{\subsubsection}
  {\fiqfuente\normalsize\bfseries\boldmath\color{fiqgris}}
  {\thesubsubsection}{0.7em}{}
\titlespacing*{\section}{0pt}{3.2ex plus 1ex minus .2ex}{1.4ex plus .2ex}
\titlespacing*{\subsection}{0pt}{2.8ex plus 1ex minus .2ex}{1.1ex plus .2ex}

% ---------- índice general ------------------------------------------
\usepackage{titletoc}
\renewcommand{\@pnumwidth}{2.2em}   % lugar para páginas de tres cifras
\titlecontents{chapter}[0pt]
  {\addvspace{11pt}\fiqfuente\bfseries\color{fiqtexto}}
  {\makebox[2em][l]{\textcolor{fiqrojo}{\thecontentslabel}}}
  {}
  {\hfill\contentspage}
  [\addvspace{2pt}]
\titlecontents{section}[2em]
  {\addvspace{1pt}}
  {\makebox[2.4em][l]{\color{fiqgris}\thecontentslabel}}
  {}
  {\titlerule*[0.7pc]{\color{fiqlinea}.}\contentspage}
\titlecontents{subsection}[4.4em]
  {\small}
  {\makebox[3em][l]{\color{fiqgris}\thecontentslabel}}
  {}
  {\titlerule*[0.7pc]{\color{fiqlinea}.}\contentspage}

% ---------- encabezados y pies --------------------------------------
\usepackage{fancyhdr}
\pagestyle{fancy}
\fancyhf{}
\newcommand{\fiqencab}{\fiqfuente\small\color{fiqgris}}
\newcommand{\fiqpagina}{{\fiqfuente\small\bfseries\color{fiqrojo}\thepage}}
\fancyhead[LE]{\fiqpagina\hspace{1em}{\fiqencab\nouppercase{\leftmark}}}
\fancyhead[RO]{{\fiqencab\nouppercase{\rightmark}}\hspace{1em}\fiqpagina}
\renewcommand{\headrulewidth}{0.4pt}
\renewcommand{\headrule}{{\color{fiqlinea}%
  \hrule height\headrulewidth width\headwidth \vskip-\headrulewidth}}
\renewcommand{\chaptermark}[1]{\markboth{%
  \ifnum\c@secnumdepth>\m@ne\ifnum\c@chapter>\z@ \thechapter\hspace{0.6em}\fi\fi #1}{}}
\renewcommand{\sectionmark}[1]{\markright{\thesection\hspace{0.6em}#1}}
\fancypagestyle{plain}{%
  \fancyhf{}\renewcommand{\headrulewidth}{0pt}%
  \fancyfoot[C]{\fiqpagina}}

% ---------- teoremas -------------------------------------------------
% Se redefinen los tres estilos de amsthm.  El color del encabezado lo
% fija la caja que rodea al entorno (\fiqcolorteo).
\newcommand{\fiqcolorteo}{fiqtexto}
\newtheoremstyle{plain}
  {8pt plus 2pt}{8pt plus 2pt}{\normalfont}{0pt}
  {\fiqfuente\bfseries\boldmath\color{\fiqcolorteo}}{.}{0.5em}
  {\thmname{#1}\thmnumber{ #2}\thmnote{\ {\mdseries\color{fiqtexto}(#3)}}}
\newtheoremstyle{definition}
  {8pt plus 2pt}{8pt plus 2pt}{\normalfont}{0pt}
  {\fiqfuente\bfseries\boldmath\color{\fiqcolorteo}}{.}{0.5em}
  {\thmname{#1}\thmnumber{ #2}\thmnote{\ {\mdseries\color{fiqtexto}(#3)}}}
\newtheoremstyle{remark}
  {8pt plus 2pt}{8pt plus 2pt}{\normalfont}{0pt}
  {\fiqfuente\bfseries\color{fiqgris}}{.}{0.5em}
  {\thmname{#1}\thmnumber{ #2}\thmnote{\ {\mdseries\color{fiqtexto}(#3)}}}

% Demostración: encabezado en Lato gris, cuadradito rojo al final.
\renewenvironment{proof}[1][\proofname]{\par
  \pushQED{\qed}\normalfont\topsep6\p@\@plus6\p@\relax
  \trivlist\item[\hskip\labelsep{\fiqfuente\bfseries\color{fiqgris}#1.}]%
  \ignorespaces}{\popQED\endtrivlist\@endpefalse}
\renewcommand{\qedsymbol}{\textcolor{fiqrojo}{\rule{0.45em}{0.45em}}}

% ---------- cajas ----------------------------------------------------
\usepackage[most]{tcolorbox}
\tcbset{
  fiqbase/.style={enhanced,breakable,frame hidden,arc=0pt,outer arc=0pt,
    colback=fiqclaro,boxrule=0pt,left=11pt,right=9pt,top=2pt,bottom=4pt,
    before skip=11pt plus 2pt,after skip=11pt plus 2pt,
    before upper={\parskip=0pt}},
  % teoremas, proposiciones, lemas, corolarios: filete rojo
  fiqteo/.style={fiqbase,
    borderline west={2.4pt}{0pt}{fiqrojo},
    before upper app={\renewcommand{\fiqcolorteo}{fiqrojooscuro}}},
  % definiciones: filete gris
  fiqdef/.style={fiqbase,
    borderline west={2.4pt}{0pt}{fiqgris}},
}
\newcommand{\fiqcajateo}[1]{\tcolorboxenvironment{#1}{fiqteo}}
\newcommand{\fiqcajadef}[1]{\tcolorboxenvironment{#1}{fiqdef}}
% se aplican al empezar el documento, cuando ya están declarados;
% los que se declaran dentro de un capítulo (teorema, proposicion, ...)
% se enganchan por nombre y funcionan igual.
\AtBeginDocument{%
  \fiqcajateo{theorem}\fiqcajateo{lemma}\fiqcajateo{corollary}%
  \fiqcajateo{proposition}\fiqcajateo{thm}\fiqcajateo{prop}%
  \fiqcajateo{coro}\fiqcajateo{lema}%
  \fiqcajateo{teorema}\fiqcajateo{proposicion}\fiqcajateo{corolario}%
  \fiqcajadef{definition}\fiqcajadef{defn}\fiqcajadef{definicion}%
}

% Ideas clave (mdframed, para que sigan valiendo las opciones que usa
% I9_pca.tex, como frametitle=...)
\usepackage{mdframed}
\renewmdenv[linewidth=0.6pt,linecolor=fiqlinea,backgroundcolor=fiqclaro,
  roundcorner=4pt,innertopmargin=9pt,innerbottommargin=9pt,
  innerleftmargin=12pt,innerrightmargin=12pt,
  frametitlefont=\fiqfuente\bfseries\color{fiqtexto},
  skipabove=10pt,skipbelow=10pt]{ideas}
% Recuadro para algoritmos
\renewmdenv[linewidth=0.8pt,linecolor=fiqgris,backgroundcolor=white,
  roundcorner=3pt,innertopmargin=7pt,innerbottommargin=7pt,
  innerleftmargin=10pt,innerrightmargin=10pt]{caja}

% Ejercicios de desarrollo (para entregar)
\newcounter{entrega}[chapter]
\renewcommand{\theentrega}{\thechapter.\arabic{entrega}}
\newenvironment{entrega}[1]{%
  \refstepcounter{entrega}%
  \begin{tcolorbox}[enhanced,breakable,colback=white,colframe=fiqrojo,
    boxrule=0.8pt,arc=3pt,left=10pt,right=10pt,top=6pt,bottom=7pt,
    colbacktitle=fiqrojo,coltitle=white,toptitle=3pt,bottomtitle=3pt,
    fonttitle=\fiqfuente\bfseries,
    before skip=14pt plus 2pt,after skip=14pt plus 2pt,
    title={Ejercicio de desarrollo \theentrega\hspace{0.6em}{\mdseries|}\hspace{0.6em}#1}]%
  \ignorespaces
}{\end{tcolorbox}}

% Fórmulas recuadradas con \boxed
\renewcommand{\boxed}[1]{{\fboxsep=5pt\fcolorbox{fiqrojo!60}{fiqrojo!4}{\m@th$\displaystyle#1$}}}

% ---------- epígrafes de figuras ------------------------------------
\usepackage{caption}
\DeclareCaptionFont{fiq}{\fiqfuente\bfseries\color{fiqgris}}
\captionsetup{font=small,labelfont=fiq,labelsep=period}

% ---------- portada --------------------------------------------------
%  En main.tex:  \title{..} \subtitulo{..} \carreras{..} \author{..}
%  \date{..} y después \maketitle.
\newcommand{\subtitulo}[1]{\def\fiq@subtitulo{#1}}
\newcommand{\carreras}[1]{\def\fiq@carreras{#1}}
\subtitulo{}\carreras{}
\renewcommand{\maketitle}{%
  \begin{titlepage}
  \thispagestyle{empty}%
  \noindent\begin{tikzpicture}[remember picture,overlay]
    % motivo circular de la marca, recortado por el borde de la hoja
    \fill[fiqrojo] ([xshift=-2.4cm,yshift=4.6cm]current page.south east) circle (5cm);
    \draw[fiqlinea,line width=0.8pt]
      ([xshift=-6.2cm,yshift=1.9cm]current page.south east) circle (2.9cm);
    \node[anchor=north west,inner sep=0pt]
      at ([xshift=2.5cm,yshift=-2.3cm]current page.north west)
      {\includegraphics[width=7.2cm]{\VA figuras/logo-fiq-color}};
  \end{tikzpicture}\par
  \vspace*{6.2cm}
  \noindent{\fiqfuente\small\color{fiqgris}\espaciado{\MakeUppercase{\fiq@subtitulo}}}\par
  \vspace{10pt}
  \noindent{\fiqfuente\bfseries\color{fiqtexto}\fontsize{44}{50}\selectfont\@title}\par
  \vspace{16pt}
  \noindent{\color{fiqrojo}\rule{2.8cm}{2.2pt}}\par
  \vspace{18pt}
  \noindent\begin{minipage}{0.62\textwidth}\raggedright
    \fiqfuente\large\color{fiqgris}\fiq@carreras
  \end{minipage}\par
  \vfill
  \noindent\begin{minipage}{0.5\textwidth}\raggedright
    {\fiqfuente\color{fiqtexto}\@author}\par\vspace{8pt}
    {\fiqfuente\small\color{fiqgris}\@date}
  \end{minipage}
  \vspace*{0.2cm}
  \end{titlepage}}

\makeatother
